
% --- Idioma y Codificación ---
\usepackage[utf8]{inputenc} % Soporte para caracteres especiales (tildes, ñ)

\usepackage[spanish,es-tabla]{babel} % Traducción de títulos (Índice, Capítulo, etc.)
\usepackage{lmodern} % Carga la fuente Latin Modern
\usepackage[T1]{fontenc} % Mejora el renderizado de caracteres
\usepackage{enumitem}
% --- Paquetes de Matemáticas (Fundamentales) ---
\usepackage{amsmath}  % Mejoras para ecuaciones y fórmulas

\usepackage{amsfonts} % Fuentes matemáticas (como el estilo blackboard bold para R, N, Z)
\usepackage{amssymb}  % Símbolos matemáticos adicionales
\usepackage{amsthm}   % Entornos para Teoremas, Lemas, Demostraciones, etc.
\usepackage{import}
\usepackage{subfiles}
% --- Otros Paquetes Útiles ---
\usepackage{graphicx}  % Para insertar imágenes y gráficas
\usepackage{xcolor}    % Para usar colores en el texto
\usepackage{hyperref}
\hypersetup{
    colorlinks=true,       % Cambia los cuadros feos por texto coloreado
    linkcolor=primaryColor, % Color para el índice y referencias internas (\ref)
    citecolor=accentColor,  % Color para las citas bibliográficas (\cite)
    urlcolor=secondaryColor % Color para los enlaces web (\url)
}
\usepackage[many]{tcolorbox}
% --- PAQUETE PARA PERSONALIZAR TÍTULOS ---
\usepackage{titlesec}

% Paquetes gráficos
\usepackage[margin=2cm]{geometry}
\usepackage{tikz}
\usetikzlibrary{shapes.geometric, calc, positioning, arrows.meta, fit, backgrounds,babel}
\usepackage{float}

\usepackage{listings} % Paquete indispensable para formatear código
\usepackage{booktabs} % Para tablas con estilo profesional
